\section{GGF Use Cases in the Benchmark}
\label{sec:ggf-use-cases}

This section documents the graph-valued functions (GGF) used in the
\texttt{xp-ggf-script} benchmark. For each operator, we summarize
\emph{(i)} its SPARQL call signature,
\emph{(ii)} the shape of the named graph it returns, and
\emph{(iii)} a compact output sample.

Unless stated otherwise, the benchmark uses the local MetaQA knowledge graph
preloaded through the CLI \texttt{--load} option, so entity-centric operators
resolve their source graph implicitly. In all cases, the operator returns a
\emph{named graph URI}, typically consumed with the following pattern:

\begin{verbatim}
BIND(ggf:OPERATOR(...) AS ?gOut)
GRAPH ?gOut {
  ...
}
\end{verbatim}

The output fragments shown below are \emph{schematic but faithful} to the real
graph shapes produced by the implementation.

\subsection{Benchmark Coverage}
\label{subsec:ggf-benchmark-coverage}

\begin{center}
\small
\begin{tabular}{l|l}
\hline
Benchmark case & GGF operator pipeline \\
\hline
\texttt{expand} & \texttt{EXPAND} \\
\texttt{paths} & \texttt{PATHS} \\
\texttt{simrank} & \texttt{SIMRANK} \\
\texttt{random\_sample} & \texttt{RANDOM-SAMPLE} \\
\texttt{local\_schema} & \texttt{LOCAL-SCHEMA} \\
\texttt{metaqa\_2hop\_local} & \texttt{EXPAND} \\
\texttt{cbd\_esbm\_summary} & \texttt{CBD} $\rightarrow$ \texttt{SUMMARY} \\
\texttt{entity\_similarity\_compare} & \texttt{CBD} $\rightarrow$ \texttt{SUMMARY} $\rightarrow$ \texttt{COMPARE-GRAPHS} \\
\texttt{entity\_similarity\_score\_only} & \texttt{CBD} $\rightarrow$ \texttt{SUMMARY} $\rightarrow$ \texttt{COMPARE-GRAPHS} \\
\hline
\end{tabular}
\end{center}

\subsection{\texttt{ggf:EXPAND}}
\label{subsec:ggf-expand}

\paragraph{Signature.}
\begin{verbatim}
ggf:EXPAND(entity_iri, hops=1, direction="both") -> graph URI
ggf:EXPAND(g_source, entity_iri, hops=1, direction="both") -> graph URI
\end{verbatim}

\paragraph{Used in.}
\texttt{expand} and \texttt{metaqa\_2hop\_local}.

\paragraph{Output schema.}
\begin{itemize}
\item The returned named graph contains only RDF triples from the expanded
neighborhood.
\item There is no dedicated metadata root node.
\item The graph is therefore queried as a plain subgraph:
\texttt{GRAPH ?gOut \{ ?s ?p ?o . \}}.
\end{itemize}

\paragraph{Representative sample output.}
\begin{verbatim}
<center> <predicate_1> <neighbor_1> .
<center> <predicate_2> "literal value" .
<neighbor_1> <predicate_3> <neighbor_2> .
\end{verbatim}

\subsection{\texttt{ggf:PATHS}}
\label{subsec:ggf-paths}

\paragraph{Signature.}
\begin{verbatim}
ggf:PATHS(start_iri, goal_iri, max_depth=4, direction="out", max_paths=10) -> graph URI
ggf:PATHS(g_source, start_iri, goal_iri, max_depth=4, direction="out", max_paths=10) -> graph URI
\end{verbatim}

\paragraph{Used in.}
\texttt{paths}.

\paragraph{Output schema.}
\begin{itemize}
\item One root node:
\begin{verbatim}
?root a slm:PathSet ;
      slm:start ?start ;
      slm:goal ?goal ;
      slm:pathCount ?count ;
      slm:hasPath ?path .
\end{verbatim}
\item For each shortest path:
\begin{verbatim}
?path a slm:Path ;
      slm:rank ?pathRank ;
      slm:length ?pathLen ;
      slm:hasStep ?step .
\end{verbatim}
\item For each step:
\begin{verbatim}
?step a slm:PathStep ;
      slm:pos ?stepPos ;
      slm:src ?src ;
      slm:pred ?pred ;
      slm:dst ?dst .
\end{verbatim}
\item If several shortest paths have the same minimal length, they may all be
returned, up to \texttt{max\_paths}.
\end{itemize}

\paragraph{Representative sample output.}
\begin{verbatim}
:root a slm:PathSet ;
      slm:start <start> ;
      slm:goal <goal> ;
      slm:pathCount 2 ;
      slm:hasPath _:p1 .

_:p1 a slm:Path ;
     slm:rank 1 ;
     slm:length 2 ;
     slm:hasStep _:s1, _:s2 .

_:s1 a slm:PathStep ;
     slm:pos 1 ;
     slm:src <start> ;
     slm:pred <bridgePred> ;
     slm:dst <mid> .

_:s2 a slm:PathStep ;
     slm:pos 2 ;
     slm:src <mid> ;
     slm:pred <goalPred> ;
     slm:dst <goal> .
\end{verbatim}

\subsection{\texttt{ggf:SIMRANK}}
\label{subsec:ggf-simrank}

\paragraph{Signature.}
\begin{verbatim}
ggf:SIMRANK(center_iri, decay=0.8, max_iter=5, top_k=10, max_nodes=200) -> graph URI
ggf:SIMRANK(g_source, center_iri, decay=0.8, max_iter=5, top_k=10, max_nodes=200) -> graph URI
\end{verbatim}

\paragraph{Used in.}
\texttt{simrank}.

\paragraph{Output schema.}
\begin{itemize}
\item One metadata root:
\begin{verbatim}
?root a cand:SimRankResult ;
      cand:center ?center ;
      cand:decay ?decay ;
      cand:iterations ?iters ;
      cand:candidate ?node .
\end{verbatim}
\item Ranked candidates:
\begin{verbatim}
?node a cand:SimRankNode ;
      cand:entity ?entity ;
      cand:score ?score ;
      cand:rank ?rank .
\end{verbatim}
\end{itemize}

\paragraph{Representative sample output.}
\begin{verbatim}
:root a cand:SimRankResult ;
      cand:center <center> ;
      cand:decay 0.8 ;
      cand:iterations 5 ;
      cand:candidate _:n1 .

_:n1 a cand:SimRankNode ;
     cand:entity <candidate_1> ;
     cand:score 0.734521 ;
     cand:rank 1 .
\end{verbatim}

\subsection{\texttt{ggf:RANDOM-SAMPLE}}
\label{subsec:ggf-random-sample}

\paragraph{Signature.}
\begin{verbatim}
ggf:RANDOM-SAMPLE(entity_iri, hops=1, direction="both", sample_rate=0.5, seed=42) -> graph URI
ggf:RANDOM-SAMPLE(g_source, entity_iri, hops=1, direction="both", sample_rate=0.5, seed=42) -> graph URI
\end{verbatim}

\paragraph{Used in.}
\texttt{random\_sample}.

\paragraph{Output schema.}
\begin{itemize}
\item The output graph contains sampled neighborhood triples \emph{and} metadata.
\item Global sample metadata:
\begin{verbatim}
?root a cand:RandomSample ;
      cand:center ?entity ;
      cand:hops ?hops ;
      cand:direction ?direction ;
      cand:sampleRate ?sampleRate ;
      cand:seed ?seed ;
      cand:exactTripleCount ?exactTripleCount ;
      cand:sampledTripleCount ?sampledTripleCount ;
      cand:estimatedTripleCount ?estimatedTripleCount ;
      cand:absoluteError ?absoluteError ;
      cand:relativeError ?relativeError ;
      cand:predicateStat ?stat .
\end{verbatim}
\item Per-predicate estimates:
\begin{verbatim}
?stat a cand:PredicateEstimate ;
      cand:rank ?rank ;
      cand:property ?property ;
      cand:exactFrequency ?exactFrequency ;
      cand:sampledFrequency ?sampledFrequency ;
      cand:estimatedFrequency ?estimatedFrequency ;
      cand:absoluteError ?predicateAbsoluteError ;
      cand:relativeError ?predicateRelativeError .
\end{verbatim}
\end{itemize}

\paragraph{Representative sample output.}
\begin{verbatim}
<center> <predicate_1> <neighbor_1> .
<center> <predicate_2> "literal value" .

:root a cand:RandomSample ;
      cand:center <center> ;
      cand:hops 1 ;
      cand:direction "both" ;
      cand:sampleRate 0.5 ;
      cand:seed 42 ;
      cand:exactTripleCount 24 ;
      cand:sampledTripleCount 12 ;
      cand:estimatedTripleCount 24.0 ;
      cand:predicateStat _:ps1 .

_:ps1 a cand:PredicateEstimate ;
      cand:rank 1 ;
      cand:property <predicate_1> ;
      cand:exactFrequency 6 ;
      cand:sampledFrequency 3 ;
      cand:estimatedFrequency 6.0 .
\end{verbatim}

\subsection{\texttt{ggf:LOCAL-SCHEMA}}
\label{subsec:ggf-local-schema}

\paragraph{Signature.}
\begin{verbatim}
ggf:LOCAL-SCHEMA(entity_iri, hops=1, direction="both", max_examples=3) -> graph URI
ggf:LOCAL-SCHEMA(g_source, entity_iri, hops=1, direction="both", max_examples=3) -> graph URI
\end{verbatim}

\paragraph{Used in.}
\texttt{local\_schema}.

\paragraph{Output schema.}
\begin{itemize}
\item Root metadata:
\begin{verbatim}
?root a cand:LocalSchema ;
      cand:center ?center ;
      cand:centerLabel ?centerLabel ;
      cand:maxHops ?maxHops ;
      cand:direction ?direction ;
      cand:slotCount ?slotCount ;
      cand:slot ?slot .
\end{verbatim}
\item Ranked schema slots:
\begin{verbatim}
?slot a cand:SchemaSlot ;
      cand:rank ?rank ;
      cand:hopCount ?hopCount ;
      cand:direction ?direction ;
      cand:property ?property ;
      cand:propertyLabel ?propertyLabel ;
      cand:frequency ?frequency ;
      cand:exampleNeighbor ?exampleNeighbor ;
      cand:exampleLiteral ?exampleLiteral .
\end{verbatim}
\item When an example neighbor has a label, the output graph also contains:
\begin{verbatim}
?exampleNeighbor rdfs:label ?exampleNeighborLabel .
\end{verbatim}
\end{itemize}

\paragraph{Representative sample output.}
\begin{verbatim}
:root a cand:LocalSchema ;
      cand:center <center> ;
      cand:maxHops 2 ;
      cand:direction "both" ;
      cand:slotCount 7 ;
      cand:slot _:slot1 .

_:slot1 a cand:SchemaSlot ;
        cand:rank 1 ;
        cand:hopCount 1 ;
        cand:direction "out" ;
        cand:property <predicate_1> ;
        cand:frequency 3 ;
        cand:exampleNeighbor <neighbor_1> ;
        cand:exampleLiteral "example literal" .

<neighbor_1> rdfs:label "Neighbor label" .
\end{verbatim}

\subsection{\texttt{ggf:CBD}}
\label{subsec:ggf-cbd}

\paragraph{Signature.}
\begin{verbatim}
ggf:CBD(entity_iri, lang="en") -> graph URI
\end{verbatim}

\paragraph{Used in.}
\texttt{cbd\_esbm\_summary}, \texttt{entity\_similarity\_compare}, and
\texttt{entity\_similarity\_score\_only}.

\paragraph{Output schema.}
\begin{itemize}
\item In the benchmark setting (local KB already loaded), the returned graph is
a plain concise bounded description around the entity:
outgoing triples plus recursive blank-node closure.
\item In local mode there is no dedicated metadata node.
\item In remote Wikidata fallback mode, an extra metadata node may be attached,
but this branch is not the one used in the MetaQA benchmark.
\end{itemize}

\paragraph{Representative sample output.}
\begin{verbatim}
<entity> <predicate_1> <neighbor_1> .
<entity> <predicate_2> "literal value" .
_:b0 <nestedPredicate> <nestedNode> .
\end{verbatim}

\subsection{\texttt{ggf:SUMMARY}}
\label{subsec:ggf-summary}

\paragraph{Signature.}
\begin{verbatim}
ggf:SUMMARY(g_source, entity_iri, k=10, neighborhood="out", mode="degree", seed=42) -> graph URI
\end{verbatim}

\paragraph{Used in.}
\texttt{cbd\_esbm\_summary}, \texttt{entity\_similarity\_compare}, and
\texttt{entity\_similarity\_score\_only}.

\paragraph{Output schema.}
\begin{itemize}
\item The returned graph contains only the selected summary triples.
\item There is no dedicated metadata root node.
\item The source graph is explicit, which makes this operator easy to compose
with \texttt{CBD} or any previously materialized graph.
\end{itemize}

\paragraph{Representative sample output.}
\begin{verbatim}
<entity> <summaryPredicate_1> <neighbor_1> .
<entity> <summaryPredicate_2> "literal value" .
<entity> <summaryPredicate_3> <neighbor_2> .
\end{verbatim}

\subsection{\texttt{ggf:COMPARE-GRAPHS}}
\label{subsec:ggf-compare-graphs}

\paragraph{Signature.}
\begin{verbatim}
ggf:COMPARE-GRAPHS(g1, g2, e1, e2) -> graph URI
\end{verbatim}

\paragraph{Used in.}
\texttt{entity\_similarity\_compare} and
\texttt{entity\_similarity\_score\_only}.

\paragraph{Output schema.}
\begin{itemize}
\item One comparison root:
\begin{verbatim}
?root a cand:Comparison ;
      cand:entity1 ?e1 ;
      cand:entity2 ?e2 ;
      cand:graph1 ?g1 ;
      cand:graph2 ?g2 ;
      cand:sharedCount ?shared ;
      cand:differingCount ?differing ;
      cand:exclusiveCount ?exclusive ;
      cand:similarityScore ?score .
\end{verbatim}
\item Property-level comparison nodes:
\begin{verbatim}
?pnode a cand:SharedProperty | cand:DifferingProperty | cand:ExclusiveProperty ;
       cand:property ?property ;
       cand:value1 ?valueToken1 ;
       cand:value2 ?valueToken2 ;
       cand:valueOverlap ?bool ;
       cand:exclusiveOf "entity1" | "entity2" .
\end{verbatim}
\item The benchmark queries currently read only the root similarity metrics, but
the output graph also retains the property-level explanation nodes.
\item These property nodes coexist in the named graph; they are not currently
linked from the root by a dedicated predicate.
\end{itemize}

\paragraph{Representative sample output.}
\begin{verbatim}
:root a cand:Comparison ;
      cand:entity1 <entity_1> ;
      cand:entity2 <entity_2> ;
      cand:graph1 <g1> ;
      cand:graph2 <g2> ;
      cand:sharedCount 4 ;
      cand:differingCount 2 ;
      cand:exclusiveCount 1 ;
      cand:similarityScore 0.571429 .

_:p1 a cand:SharedProperty ;
     cand:property <predicate_1> ;
     cand:value1 "token_a" ;
     cand:value2 "token_a" ;
     cand:valueOverlap true .

_:p2 a cand:ExclusiveProperty ;
     cand:property <predicate_2> ;
     cand:value1 "token_b" ;
     cand:exclusiveOf "entity1" .
\end{verbatim}

\subsection{Composed Use Cases}
\label{subsec:ggf-composed-use-cases}

Some benchmark cases are direct one-operator calls, while others are explicit
server-side compositions of graph-valued functions.

\paragraph{\texttt{cbd\_esbm\_summary}.}
For each candidate entity, the benchmark computes:
\begin{verbatim}
?gCBD     := ggf:CBD(?entity)
?gSummary := ggf:SUMMARY(?gCBD, ?entity, 10, "both", "degree", 42)
\end{verbatim}
The final query then reads only the selected triples from \texttt{?gSummary}.

\paragraph{\texttt{entity\_similarity\_compare}.}
For a reference entity and each candidate, the benchmark computes:
\begin{verbatim}
?gRefCBD     := ggf:CBD(?refEntity)
?gRefSummary := ggf:SUMMARY(?gRefCBD, ?refEntity, ...)
?gEntCBD     := ggf:CBD(?entity)
?gEntSummary := ggf:SUMMARY(?gEntCBD, ?entity, ...)
?gCmp        := ggf:COMPARE-GRAPHS(?gRefSummary, ?gEntSummary, ?refEntity, ?entity)
\end{verbatim}
The final result projects only the root-level similarity statistics.

\paragraph{\texttt{entity\_similarity\_score\_only}.}
This case reuses the same pipeline as
\texttt{entity\_similarity\_compare}, but projects only
\texttt{cand:similarityScore}.

\paragraph{\texttt{metaqa\_2hop\_local}.}
This case batches multiple questions in a single SPARQL query. For each anchor
entity, it materializes a 2-hop expanded graph with \texttt{ggf:EXPAND(...)}
and resolves answer candidates by pattern matching directly inside that
temporary graph.

\subsection{Practical Reading Guide}
\label{subsec:ggf-reading-guide}

The benchmark intentionally mixes two styles of output graphs:
\begin{itemize}
\item \textbf{Plain triple graphs:} \texttt{EXPAND}, \texttt{CBD}, and
\texttt{SUMMARY} return ordinary RDF triples without an explicit metadata root.
\item \textbf{Structured result graphs:} \texttt{PATHS}, \texttt{SIMRANK},
\texttt{RANDOM-SAMPLE}, \texttt{LOCAL-SCHEMA}, and
\texttt{COMPARE-GRAPHS} return typed result nodes carrying metadata and
ranked substructures.
\end{itemize}

This distinction matters for SPARQL authoring. Plain triple graphs are best
queried with generic triple patterns, whereas structured result graphs should
usually be queried through their root node type and typed child nodes.
